\section{Преобразование Абеля (суммирование по частям)}

Тонкие признаки сходимости позволяют исследовать ряды $\sum a_k b_k$, к которым неприменимы ни
признак Лейбница, ни признаки сравнения. В их основе лежит \emph{преобразование Абеля} —
дискретный аналог интегрирования по частям.

\begin{lemma}[преобразование Абеля]
Пусть $B_n=\sum_{k=1}^n b_k$, $B_0=0$. Тогда
\[
  \sum_{k=1}^n a_k b_k=a_n B_n+\sum_{k=1}^{n-1}(a_k-a_{k+1})B_k .
\]
\end{lemma}

\begin{proof}
Подставим $b_k=B_k-B_{k-1}$ и сгруппируем по $B_k$:
\[
  \sum_{k=1}^n a_k(B_k-B_{k-1})
   =\sum_{k=1}^{n-1}B_k(a_k-a_{k+1})+a_n B_n .
\]
Здесь $a_k$ играет роль «множителя», $b_k$ — «дифференциала»; разность $a_k-a_{k+1}$ заменяет
производную, а суммирование $B_k$ — первообразную. Это полностью повторяет схему
$\int u\,dv=uv-\int v\,du$.
\end{proof}

\section{Признак Дирихле}

\begin{theorem}[Дирихле]
Ряд $\displaystyle\sum_{k=1}^\infty a_k b_k$ сходится, если:
\begin{enumerate}[label=\arabic*),leftmargin=*]
  \item $a_k$ монотонно стремится к нулю ($a_k\downarrow0$);
  \item частичные суммы $B_n=\sum_{k=1}^n b_k$ ограничены в совокупности: $\abs{B_n}\le M$.
\end{enumerate}
\end{theorem}

\begin{proof}
По преобразованию Абеля $\displaystyle S_n=a_n B_n+\sum_{k=1}^{n-1}(a_k-a_{k+1})B_k$.
\emph{Первое слагаемое} $a_n B_n\to0$, так как $a_n\to0$, а $\abs{B_n}\le M$.
\emph{Второе слагаемое} — частичная сумма абсолютно сходящегося ряда:
\[
  \sum_{k=1}^{\infty}\abs{a_k-a_{k+1}}\abs{B_k}\le M\sum_{k=1}^{\infty}(a_k-a_{k+1})
   =M\,a_1<\infty
\]
(модули сняты в силу монотонности, ряд телескопический). Значит, существует
$\displaystyle\lim_{n\to\infty}S_n$.
\end{proof}

\begin{example}[классический пример $\sum\frac{\sin n}{n^p}$]
Ряд $\displaystyle\sum_{n=1}^\infty\frac{\sin n}{n^p}$ при $p>0$ сходится. Берём
$a_n=\dfrac1{n^p}\downarrow0$ и $b_n=\sin n$; частичные суммы синусов ограничены:
\[
  B_n=\sum_{k=1}^n\sin k=\frac{\cos\frac12-\cos\bigl(n+\frac12\bigr)}{2\sin\frac12},\qquad
  \abs{B_n}\le\frac{1}{\sin\frac12}.
\]
По признаку Дирихле ряд сходится. При $0<p\le1$ сходимость условная, при $p>1$ — абсолютная.
\end{example}

\section{Признак Абеля}

\begin{theorem}[Абель]
Ряд $\displaystyle\sum_{k=1}^\infty a_k b_k$ сходится, если:
\begin{enumerate}[label=\arabic*),leftmargin=*]
  \item $a_k$ монотонна и ограничена;
  \item ряд $\sum b_k$ сходится.
\end{enumerate}
\end{theorem}

\begin{proof}
Монотонная ограниченная последовательность $a_k$ имеет конечный предел $a$. Частичные суммы $B_n$
сходятся (ряд $\sum b_k$ сходится), значит ограничены: $\abs{B_k}\le M$. Снова по преобразованию
Абеля $S_n=a_n B_n+\sum_{k=1}^{n-1}(a_k-a_{k+1})B_k$: первое слагаемое имеет предел $a\cdot B$
(где $B=\sum b_k$), второе сходится абсолютно, как в признаке Дирихле. Поэтому $\lim S_n$
существует.
\end{proof}

\begin{remark}[различие условий]
В обоих признаках $a_k$ монотонна. Различие — в требованиях:
\begin{itemize}[leftmargin=*]
  \item \textbf{Дирихле}: $a_k\to0$ (сильное условие на $a_k$), но от $\sum b_k$ требуется лишь
        \emph{ограниченность} частичных сумм (ряд может расходиться, как $\sum\sin k$);
  \item \textbf{Абель}: $a_k$ лишь \emph{ограничена} (предел может быть ненулевым), зато $\sum b_k$
        обязан \emph{сходиться}.
\end{itemize}
Признак Абеля сводится к Дирихле: если $a_k\to a$, то $a_k-a\to0$ монотонно, а
$\sum(a_k-a)b_k$ исследуется по Дирихле, $\sum a\,b_k$ сходится по условию.
\end{remark}

\begin{example}
$\displaystyle\sum_{k=1}^\infty\frac{(-1)^k}{\sqrt k}\Bigl(1+\frac1k\Bigr)^k$: множитель
$a_k=\bigl(1+\tfrac1k\bigr)^k$ возрастает и ограничен ($0\le a_k<e$), а ряд
$\sum\dfrac{(-1)^k}{\sqrt k}$ сходится по признаку Лейбница. По признаку Абеля исходный ряд
сходится.
\end{example}
